\centerline{\bf \Huge 8 класс}

\problem{\textbf{1}} \textbf{Критерии:} \\
\textit{7 баллов} $-$ Любое полное верное решение с найденным $x=256$, с доказательством существования ещё двух других действительных корней и их примерным нахождением $x=-0,98$ ;$x=1,02$.\\
\textit{6 баллов} $-$ Любое полное верное решение с найденным $x=256$ и с доказательством существования ещё двух других действительных корней.\\
\textit{5 баллов} $-$ Найден $x=256$ с обоснованием или с проверкой.\\
\textit{2 балла} $-$ Найден $x=256$ без обоснований и без проверки.\\
\textit{0-2 балла} $-$ Верный $x=256$ не найден, но имеются математически верные продвижения.\\
\textit{0 баллов} $-$ Любое неверное решение или неверный ответ.\\
\textit{0 баллов} $-$ Решение отсутствует.\\
\textit{0 баллов} $-$ При обоснованных подозрениях на списывание участником из любого источника (Интернет, учебное пособие, ключи олимпиады и т.д.) член жюри имеет право указать на такие фрагменты с помощью восклицательных знаков на полях. \\

\problem{\textbf{2}} \textbf{Критерии суммируются:} \\
\textit{0-4 балла} $-$ Решение А.\\
\textit{0-1 балл} $-$ Составление текстового условия Б.\\
\textit{0-2 балла} $-$ Любое правильное решение Б.\\
\textit{0 баллов} $-$ Любой ответ без объяснений.\\
\textit{0 баллов} $-$ Решение отсутствует.\\
\textit{0 баллов} $-$ При обоснованных подозрениях на списывание участником из любого источника (Интернет, учебное пособие, ключи олимпиады и т.д.) член жюри имеет право указать на такие фрагменты с помощью восклицательных знаков на полях. \\



\problem{\textbf{3}} \textbf{Критерии суммируются:} \\
\textit{2 балла} $-$ Один слон половину дня пьёт воду "из озера", а половину - "из ключей". Поэтому ему понадобится $182,5 \cdot 2=365$ дней.\\
\textit{2 балла} $-$ Доказано, что в озере (без ключей) 182,5 дневных порций слона.\\
\textit{3 балла} $-$ Доказано, что ключи восполняют за день половину дневной порции слона.\\
\textit{0 баллов} $-$ Любой ответ без объяснений.\\
\textit{0 баллов} $-$ Решение неверное, продвижения отсутствуют.\\
\textit{0 баллов} $-$ Решение отсутствует.\\
\textit{0 баллов} $-$ При обоснованных подозрениях на списывание участником из любого источника (Интернет, учебное пособие, ключи олимпиады и т.д.) член жюри имеет право указать на такие фрагменты с помощью восклицательных знаков на полях. \\

\textbf{ИЛИ:}
\textit{7 баллов} $-$ Другое полное верное решение.\\



\problem{\textbf{4}} \textbf{Критерии суммируются:} \\
\textit{2 балла} $-$ Доказано, что трапеция обязана быть равнобокой.\\
\textit{1 балл} $-$ Указано, что трапеция обязана быть равнобокой, но не доказано.\\
\textit{1 балл} $-$ Так как площадь 32, то высота трапеции равна $32:8=4$.\\
\textit{1 балл} $-$ Средняя линия трапеции равна полусумме оснований трапеции.\\
\textit{0 баллов} $-$ Любой ответ без объяснений.\\
\textit{0 баллов} $-$ Решение неверное, продвижения отсутствуют.\\
\textit{0 баллов} $-$ Решение отсутствует.\\
\textit{0 баллов} $-$ При обоснованных подозрениях на списывание участником из любого источника (Интернет, учебное пособие, ключи олимпиады и т.д.) член жюри имеет право указать на такие фрагменты с помощью восклицательных знаков на полях.\\

\textbf{ИЛИ:}
\textit{7 баллов} $-$ Полное верное решение.\\


\problem{\textbf{5}} \textbf{Критерии:} \\
\textit{7 баллов} $-$ Полное верное решение.\\
\textit{4-6 баллов} $-$ Составлена правильная система уравнений, но до конца не решена.\\
\textit{0-1 балла} $-$ Оценка стороны квадрата.\\
\textit{0 баллов} $-$ Любой ответ без объяснений.\\
\textit{0 баллов} $-$ Решение неверное, продвижения отсутствуют.\\
\textit{0 баллов} $-$ Решение отсутствует.\\
\textit{0 баллов} $-$ При обоснованных подозрениях на списывание участником из любого источника (Интернет, учебное пособие, ключи олимпиады и т.д.) член жюри имеет право указать на такие фрагменты с помощью восклицательных знаков на полях. \\